\documentclass[journal,twocolumn,10pt]{IEEEtran}

\usepackage{cite}
\usepackage{amsmath,amssymb,amsfonts}
\usepackage{algorithmic}
\usepackage{graphicx}
\usepackage{textcomp}
\usepackage{xcolor}
\usepackage{booktabs}
\usepackage{multirow}
\usepackage{array}
\usepackage{url}

\graphicspath{{../figures/vector_pdf/}{../figures/}}

\begin{document}

\title{MARLIN-Twin: Maritime Adaptive Resilience Learning with Integrated Networked Twins for Autonomous Vessel Collision Avoidance Under Communication Outages}

\author{Deugs~Research~Group
\thanks{Manuscript received July 29, 2026. This work was supported by the Maritime Autonomous Systems Research Initiative.}
}

\markboth{IEEE Transactions on Intelligent Transportation Systems}%
{MARLIN-Twin: Maritime Adaptive Resilience Learning with Integrated Networked Twins}

\maketitle

\begin{abstract}
Cooperative collision avoidance among Autonomous Surface Vehicles (ASVs) requires reliable inter-vessel state awareness and strict COLREGs adherence despite Automatic Identification System (AIS) dropouts and Vehicle-to-Vehicle bandwidth degradation ($\lambda \in [0,1]$). We present \textbf{MARLIN-Twin}, a multi-agent reinforcement learning framework combining an onboard EKF/JPDA Digital Twin, Multi-Head Graph Attention (GAT) observation encoding, and CTDE MAPPO, together with a validated 3-DOF MMG hydrodynamic simulation testbed, and use both to evaluate -- rather than assume -- how well this design delivers degradation-resilient collision avoidance. An earlier version of this evaluation, after fixing a genuine IMO Res. MSC.137(76) turning-circle non-compliance and a missing per-seed re-initialization bug, reported Rule-Based COLREGs reproducibly outperforming every RL-trained method on this paper's 2-vessel head-on benchmark. Investigating that result surfaced a further, more consequential bug: a curriculum-training seed-collapse fault that silently replayed one episode for 40\% of every RL policy's training, alongside a Rule-Based action-convention mismatch and a training/evaluation metric inconsistency, all of which we diagnose and fix. Under the resulting corrected pipeline and a realistic 500-episode training budget (a $3.3\times$ increase over the previous 150), MARLIN-Twin significantly outperforms Rule-Based COLREGs (safety score $0.097$ vs.\ $0.085$; Welch's $t$-test $p=0.00018$, Cohen's $d=8.5$) -- reversing the earlier headline finding. We further add multi-agent Soft Actor-Critic (SAC) as a fifth algorithm to test whether this result is architecture-specific: MARLIN-Twin remains statistically indistinguishable from MADDPG, Independent PPO, and SAC, so which learned algorithm is used remains an open question even though the RL-vs-rule-based question is now resolved on this scenario. A reward-reweighting sweep across eight configurations finds the published reward weighting already performs near the best of the sweep, and that naively upweighting the safety term alone destabilizes training rather than reliably improving it. Screening five harder, more ambiguous scenarios against Rule-Based COLREGs' own confirmed structural weakness -- it collapses every neighbor to the single nearest vessel and cannot represent simultaneous conflicting COLREGs obligations -- identifies two, a converging multi-sector port approach and a restricted-visibility crossing (both $N=6$), on which every trained RL variant, not only MARLIN-Twin, clearly outperforms Rule-Based COLREGs; a fourth, cleaner multi-vessel channel-convergence scenario continues to favor Rule-Based COLREGs, consistent with it never presenting the conflicting-obligation failure mode the other two do. We validate the EKF/JPDA Digital Twin against a real NOAA MarineCadastre AIS trajectory (497m overall RMSE, 1539m during an injected 300s blackout). We report every seeding, measurement-harness, and metric-consistency bug found across both the original and this follow-up evaluation, and the statistical conclusions that changed once each was fixed, as a central part of this paper's contribution: an evaluation's comparative claims are only as trustworthy as its randomization and measurement infrastructure, and a result that reverses under closer scrutiny is worth reporting precisely because it did.
\end{abstract}

\begin{keywords}
Autonomous Surface Vehicles (ASVs), Multi-Agent Reinforcement Learning (MARL), Digital Twins, Extended Kalman Filter (EKF), Graph Attention Networks (GAT), COLREGs, Communication Resilience.
\end{keywords}

\section{Introduction}
\IEEEPARstart{A}{utonomous} Surface Vehicles (ASVs) are poised to revolutionize maritime transportation by enhancing navigational safety and enabling cooperative maneuvering \cite{zheng2025marl}. However, executing cooperative collision avoidance maneuvers in dense multi-vessel environments remains a major challenge \cite{xu2024colregs, wang2025realtime_colregs}. Vessels must negotiate complex multi-vessel encounter scenarios while strictly adhering to the International Regulations for Preventing Collisions at Sea (COLREGs) \cite{tang2023colregs}.

In practice, maritime V2V communication networks are subject to severe environmental disturbances, including multipath fading, sea spray attenuation, AIS transponder congestion, and intentional signal jamming \cite{alqurashi2024maritime}. Existing MARL algorithms (such as IPPO and MADDPG) often rely on full, uncorrupted observation sharing or continuous high-bandwidth telemetry \cite{hu2023eventtriggered}. When communication quality degrades ($\lambda \to 0$), these baseline systems suffer severe performance degradation, resulting in heightened collision risks and COLREGs violations \cite{gu2024resilient_marl}.

Toward these limitations, we present \textbf{MARLIN-Twin}, a multi-agent framework and validated simulation testbed that combines onboard EKF/JPDA Digital Twins with Graph Attention Networks (GAT) and CTDE MAPPO policy optimization, and use it to evaluate how well this architecture -- and the graph-attention/digital-twin design choices within it -- actually deliver degradation-resilient collision avoidance, rather than assume the answer. As illustrated in Fig.~\ref{fig:system_arch}, MARLIN-Twin equips each ASV with a local Digital Twin that estimates vessel states and reconstructs trajectories during AIS outages via Dead Reckoning; Section~\ref{sec:experimental-evaluation} reports where this design measurably helps and where, under the encounter geometry and training budget tested here, it does not.

\begin{figure*}[t]
\centering
\includegraphics[width=0.95\textwidth]{fig1_system_architecture_ieee.pdf}
\caption{Overall MARLIN-Twin System Architecture Block Diagram. Raw sensor inputs feed into an onboard EKF/JPDA Digital Twin, which constructs a spatial graph processed by a Multi-Head GAT policy trained via CTDE MAPPO.}
\label{fig:system_arch}
\end{figure*}

The main contributions of this work are fourfold:
\begin{enumerate}
    \item \textbf{3-DOF MMG Hydrodynamic Engine Evaluated Against IMO Sea Trial Criteria}: We implement a 3-DOF Manoeuvring Modeling Group (MMG) nonlinear vessel model and evaluate it against both IMO Resolution MSC.137(76) sea trial criteria \cite{imo2002standards}: the model meets both the 10$^\circ$/10$^\circ$ Zig-Zag overshoot criterion and the 30$^\circ$ Turning Circle criterion, the latter after diagnosing and fixing a latent measurement bug in our own zig-zag test harness that had been masking a viable coefficient fix (Section~\ref{sec:limitations} details both).
    \item \textbf{EKF/JPDA Digital Twin Outage Recovery}: We deploy an onboard state estimation engine that fuses noisy AIS/Radar signals and maintains continuous track estimates during 300s AIS blackouts.
    \item \textbf{Spatial GAT \& 128-Bit Payload Protocol}: We construct dynamic relative observation graphs utilizing Multi-Head GAT attention weights ($\alpha_{ij}$) and serialize messages into a 128-bit binary payload for low-bandwidth mesh channels.
    \item \textbf{Empirical Resilience Study, a Realistic-Budget Reversal, and a Harder-Scenario Battery, Including Every Infrastructure Bug We Found and Fixed}: An earlier version of this evaluation, after fixing a missing per-seed re-initialization bug, found MARLIN-Twin's absolute safety score statistically indistinguishable from MADDPG/Independent PPO but reproducibly below Rule-Based COLREGs. Investigating that result surfaced a further curriculum-training seed-collapse bug (Section~\ref{sec:limitations}); fixing it and retraining at a realistic 500-episode budget reverses the finding -- MARLIN-Twin now significantly outperforms Rule-Based COLREGs ($p=0.00018$, Cohen's $d=8.5$). Adding multi-agent SAC as a fifth algorithm shows this reversal is not architecture-specific (every RL variant improves), but MARLIN-Twin still cannot be distinguished from MADDPG, IPPO, or SAC. A reward-reweighting sweep and a five-scenario harder-encounter battery (Section~\ref{sec:experimental-evaluation}) further test this result's robustness, finding two additional scenarios -- exploiting Rule-Based COLREGs' confirmed single-nearest-neighbor limitation -- on which every RL variant outperforms it. We document every bug found across both evaluation passes alongside the results themselves (Section~\ref{sec:limitations}) rather than only the results that survived them.
\end{enumerate}

\section{Related Work}
\subsection{COLREGs-Compliant Motion Planning}
Traditional ASV motion planning relies on Velocity Obstacles (VO), Artificial Potential Fields (APF), and Model Predictive Control (MPC) \cite{wang2025realtime_colregs, xu2024colregs}. These methods offer deterministic safety guarantees and, as this paper's own results show (Section~\ref{sec:experimental-evaluation}), can be highly effective in a single, well-specified pairwise encounter; the case against them is scalability and flexibility as encounter density and rule ambiguity grow (computational complexity and rigid rule interpretation under high-density multi-vessel encounters), not that they are unconditionally worse than learned approaches.

\subsection{Multi-Agent Reinforcement Learning in Maritime Autonomy}
MARL algorithms have emerged as powerful tools for autonomous navigation \cite{zheng2025marl, yu2023mappo}. Independent PPO (IPPO) treats other vessels as static environmental components, leading to non-stationarity. Multi-Agent DDPG (MADDPG) utilizes a centralized critic but requires complete state availability \cite{lowe2017multi}. MARLIN-Twin investigates whether GAT observation aggregation and CTDE MAPPO \cite{velickovic2018graph, park2025gat} can address these limitations; Section~\ref{sec:experimental-evaluation} reports the extent to which they do.

\subsection{Digital Twins and State Estimation under Outages}
Digital Twin technology links physical assets with digital representations \cite{chen2024digitaltwin, chen2025digital_twin_fleet, zhang2025digital_twin_asv}. In maritime domain, EKF and JPDA track association enable robust multi-target state estimation \cite{zhang2023ekf}. MARLIN-Twin extends these techniques to maintain agent policy stability during prolonged blackout periods \cite{wu2024ais_outage, huang2025ais, zhou2024ais_denoising}.

\section{3-DOF Hydrodynamic Vessel Model \& Problem Formulation}
\subsection{3-DOF MMG Vessel Dynamics}
\label{sec:mmg-dynamics}
The 3-DOF horizontal manoeuvring motion of an ASV (surge $u$, sway $v$, yaw rate $r$) is governed by the MMG equations \cite{yoshimura2009mmg, fossen2011handbook, yasukawa2015mmg}:
\begin{equation}
(m + m_x) \dot{u} - (m + m_y) v r = X_{\text{hydro}} + X_{\text{prop}} + X_{\text{rud}}
\end{equation}
\begin{equation}
(m + m_y) \dot{v} + (m + m_x) u r = Y_{\text{hydro}} + Y_{\text{rud}}
\end{equation}
\begin{equation}
(I_z + J_z) \dot{r} = N_{\text{hydro}} + N_{\text{rud}}
\end{equation}
where $m$ is vessel displacement, $m_x, m_y, J_z$ are added masses/moment of inertia, and $X, Y, N$ represent hydrodynamic forces and moments \cite{wang2024mmg, yasukawa2015mmg}. Kinematic position update follows $\dot{\boldsymbol{\eta}} = \mathbf{R}(\psi)\boldsymbol{\nu}$, solved via 4th-order Runge-Kutta (RK4) integration.

\begin{figure}[t]
\centering
\includegraphics[width=0.48\textwidth]{fig5_sea_trials_ieee.pdf}
\caption{3-DOF MMG Hydrodynamic Model Sea Trial Validation against IMO Resolution MSC.137(76) standards (30$^\circ$ Starboard Turning Circle and 10$^\circ$/10$^\circ$ Zig-Zag Maneuver).}
\label{fig:sea_trials}
\end{figure}

As shown in Fig. \ref{fig:sea_trials}, the model reaches a 1st overshoot angle of 10.0$^\circ$ and a 2nd overshoot angle of 19.6$^\circ$ on the 10$^\circ$/10$^\circ$ zig-zag maneuver. IMO MSC.137(76) \cite{imo2002standards} sets separate, speed-dependent limits for each: at this vessel's $L/V \approx 18.75$s (150m length, 8 m/s), the applicable limits are $5 + 0.5(L/V) \approx 14.4^\circ$ for the first overshoot and $17.5 + 0.75(L/V) \approx 31.6^\circ$ for the second, both comfortably cleared. The 30$^\circ$ turning circle -- the maximum rudder angle this vessel model supports, per Table~\ref{tab:hyperparams} -- yields a tactical diameter of 466 m over its first complete 360$^\circ$ loop, comfortably inside the IMO standard's $5L$ ceiling of 750 m for this vessel's 150 m length. Both criteria continue to hold when the zig-zag maneuver is extended past the two officially-checked reversals to 10 consecutive reversal cycles (peak overshoot 1.84$^\circ$ across the extended run), confirming a genuinely damped response rather than one that merely satisfies the two checked overshoots en route to a slower divergence (see Section~\ref{sec:limitations} for the coefficient search and a test-harness bug this check surfaced and fixed along the way).

\subsection{POMDP Game Formulation}
The multi-ASV interaction is formulated as a Partially Observable Markov Decision Process (POMDP) defined by $\langle \mathcal{N}, \mathcal{S}, \{\mathcal{A}_i\}, \{\mathcal{O}_i\}, \mathcal{P}, \mathcal{R}, \gamma \rangle$, where $\mathcal{N} = \{1, \dots, N\}$ is the set of $N$ vessels (Section~\ref{sec:experimental-evaluation} specifies the value of $N$ used in each experiment).

Each vessel $i$ receives an observation $o_i \in \mathcal{O}_i$ containing its own state $\mathbf{s}_i = [x_i, y_i, \psi_i, u_i, v_i, r_i]^T$ and perceived neighbor states $\hat{\mathbf{s}}_j$. Action $a_i \in \mathcal{A}_i = [n_{\text{rpm}}, \delta_{\text{rudder}}, p_{\text{msg}}]^T$ controls propeller RPM, rudder angle, and message priority.

Vessel $i$'s per-step reward is a fixed-weight sum of three terms,
\begin{equation}
r_i = w_{\text{safety}} \, r_i^{\text{safety}} + w_{\text{colregs}} \, r_i^{\text{colregs}} + w_{\text{eff}} \, r_i^{\text{eff}},
\label{eq:reward}
\end{equation}
with $(w_{\text{safety}}, w_{\text{colregs}}, w_{\text{eff}}) = (2.0, 1.0, 1.0)$ throughout every experiment in this paper. $r_i^{\text{safety}} = -\exp(-d^{\text{cpa}}_i / 200)$ penalizes a close current-encounter closest-point-of-approach distance $d^{\text{cpa}}_i$ (5000m, i.e. no penalty, if vessel $i$ has no active encounter that step); $r_i^{\text{colregs}}$ is the product of \texttt{COLREGsEngine.evaluate\_compliance}'s per-encounter compliance score ($\in[0,1]$, 1.0 with no active encounter) across every encounter vessel $i$ is party to that step; $r_i^{\text{eff}} = -d^{\text{goal}}_i / 5000$ penalizes remaining distance to the current waypoint ($1.0$, a bonus, once the route is fully complete). The team reward used for the value baseline is $\bar{r} = \frac{1}{N}\sum_i r_i$. We call out explicitly that $r_i^{\text{safety}}$ is a per-step, single nearest-encounter proxy distinct from the safety score $J(\lambda)$ (Eq.~\ref{eq:safety_score}, Section~\ref{sec:experimental-evaluation}) used for every comparison in this paper's results: $J(\lambda)$ is computed post-hoc from each full episode's true minimum pairwise separation and was not used as a training signal, so a policy is never directly optimized against the metric this paper evaluates it on.

\begin{figure}[t]
\centering
\includegraphics[width=0.48\textwidth]{fig2_digital_twin_flowchart_ieee.pdf}
\caption{Digital Twin EKF/JPDA State Estimation \& Outage Recovery Flowchart. Seamlessly transitions to Dead Reckoning fallback during AIS signal blackouts.}
\label{fig:flowchart}
\end{figure}

\begin{figure}[t]
\centering
\includegraphics[width=0.48\textwidth]{fig6_digital_twin_blackout_ieee.pdf}
\caption{EKF/JPDA Digital Twin Trajectory Estimation and Position Tracking Error during a 300s AIS Blackout Event.}
\label{fig:blackout_tracking}
\end{figure}

\section{The MARLIN-Twin Framework}
\label{sec:framework}
\subsection{Digital Twin EKF/JPDA Estimation Engine}
The onboard Digital Twin maintains state estimate $\hat{\mathbf{s}}_i$ and error covariance $\mathbf{P}_i$ via EKF update when AIS signals are present \cite{bar-shalom2001tracking, zhang2023ekf}. Under AIS blackout ($\lambda \to 0$), the estimator switches to Dead Reckoning prediction:
\begin{equation}
\hat{\mathbf{s}}_{k|k-1} = f(\hat{\mathbf{s}}_{k-1|k-1}, \mathbf{u}_{k-1}), \quad \mathbf{P}_{k|k-1} = \mathbf{F}_k \mathbf{P}_{k-1|k-1} \mathbf{F}_k^T + \mathbf{Q}_k
\end{equation}

Fig. \ref{fig:flowchart} and Fig. \ref{fig:blackout_tracking} illustrate the decision logic governing signal processing, track maintenance, and error propagation.

\begin{figure}[t]
\centering
\includegraphics[width=0.42\textwidth]{fig3_gat_attention_diagram_ieee.pdf}
\caption{Multi-Head GAT Graph Attention Mechanism: real attention weights $\alpha_{ij}$ (mean over $K=4$ heads, including the ownship self-loop) computed by a forward pass of the trained seed-42 checkpoint's two per-vessel encoders on a constructed five-vessel scene with one neighbor per COLREGs encounter type. Each panel's five weights sum to 1.0.}
\label{fig:gat_diagram}
\end{figure}

\begin{figure}[t]
\centering
\includegraphics[width=0.48\textwidth]{fig8_degradation_heatmap_ieee.pdf}
\caption{Inter-Vessel Bandwidth Utilization Heatmap Matrix under varying communication quality levels ($\lambda$).}
\label{fig:bandwidth_heatmap}
\end{figure}

Fig.~\ref{fig:bandwidth_heatmap} sweeps MARLIN-Twin's seed-42 checkpoint over the same six $\lambda$ levels used throughout this paper, reporting both the channel's actual bandwidth (bits/s, scaling linearly with $\lambda$ by construction of \texttt{CommunicationChannelManager.set\_degradation}) and the resulting mean team reward $\bar{r}$ (Eq.~\ref{eq:reward}). Bandwidth scales as expected; team reward is flat across every level tested. This is a distinct finding from, and should not be read as corroborating, the safety-score invariance discussed in Section~\ref{sec:limitations} (\emph{Resilience index interpretation}): $\bar{r}$'s own safety term still uses a per-step, velocity-only closest-point-of-approach projection rather than the true minimum pairwise separation $J(\lambda)$ (Eq.~\ref{eq:safety_score}) is computed from, and that reward term was deliberately left unchanged in this study (only the post-hoc evaluation metric was corrected) -- so this flat reward curve is not yet expected to reflect degradation regardless of what the corrected safety score shows.

\subsection{Multi-Head Graph Attention Network (GAT)}
The spatial relational topology of surrounding vessels is represented as a dynamic graph $\mathcal{G} = (\mathcal{V}, \mathcal{E})$. GAT attention coefficients $\alpha_{ij}$ between vessel $i$ and neighbor $j$ are computed as \cite{velickovic2018graph, li2025gat_maritime}:
\begin{equation}
\alpha_{ij} = \frac{\exp\left(\text{LeakyReLU}\left(\mathbf{a}^T [\mathbf{W}h_i \,||\, \mathbf{W}h_j]\right)\right)}{\sum_{k \in \mathcal{N}_i} \exp\left(\text{LeakyReLU}\left(\mathbf{a}^T [\mathbf{W}h_i \,||\, \mathbf{W}h_k]\right)\right)}
\end{equation}
Fig.~\ref{fig:gat_diagram} reports the actual attention weights computed by both vessels' trained encoders on this constructed scene, rather than an illustrative example (from the checkpoint retrained under the corrected turning-circle dynamics, encounter geometry, and seeding; Section~\ref{sec:limitations}). Vessel 0's encoder assigns $\alpha = 0.221$ to the head-on neighbor, $0.197$ to the give-way crossing neighbor, $0.170$ to the overtaking neighbor, and $0.191$ to the stand-on crossing neighbor, retaining $0.220$ on its own self-loop -- its highest weight, narrowly, goes to the head-on neighbor. Vessel 1's encoder, run on the identical scene, assigns $0.192$, $0.200$, $0.207$, and $0.201$ respectively, retaining $0.200$ on itself -- its highest weight goes to the overtaking neighbor, not the head-on one. The two encoders therefore disagree with each other on which encounter type receives the most attention, and neither ranking is far from uniform ($0.170$--$0.221$ across a $0.20$ mean, i.e.\ a narrow spread). As discussed in Section~\ref{sec:limitations}, we attribute this to the checkpoint's limited training (150 episodes on a 2-vessel scenario) being evaluated here out-of-distribution on a 5-vessel scene, and we do not find evidence in this checkpoint that GAT attention reliably prioritizes higher-risk encounter types \emph{consistently across vessels}; Fig.~\ref{fig:gat_diagram} is presented as an honest illustration of the attention mechanism itself rather than as a validated finding about learned risk prioritization.

\subsection{128-Bit Payload Binary Protocol}
To minimize bandwidth consumption, inter-vessel messages are serialized into a 128-bit payload:
\begin{equation}
\text{Payload}_{128} = [\text{ID}_{8} \,||\, \text{Lat}_{24} \,||\, \text{Lon}_{24} \,||\, \text{SOG}_{16} \,||\, \text{COG}_{16} \,||\, \text{Hdg}_{16} \,||\, \text{CRC}_{24}]
\end{equation}
A JSON-serialized message carrying the same six fields (e.g.\ \texttt{\{"id":12,"lat":40.7128,"lon":-74.0060,"sog":12.5,"cog":245.3,"hdg":244.8\}}, 84 ASCII characters $=$ 672 bits) is, by this direct comparison, 5.25$\times$ larger than the 128-bit payload -- an 81\% reduction in transmission footprint for this representative message; the exact percentage depends on field-name verbosity and numeric precision in the JSON encoding chosen for comparison.

\subsection{CTDE MAPPO Training Infrastructure}
We optimize joint policy $\pi_\theta$ using Centralized Training with Decentralized Execution (CTDE) MAPPO with Generalized Advantage Estimation (GAE, $\gamma=0.99, \lambda_{\text{gae}}=0.95$) \cite{yu2023mappo, schulman2017proximal}. The clipped surrogate objective for each agent $i$ is formulated as:
\begin{equation}
L^{\text{CLIP}}(\theta) = \hat{\mathbb{E}}_t \left[ \min\left(r_t(\theta)\hat{A}_t^{\text{GAE}}, \text{clip}(r_t(\theta), 1-\epsilon, 1+\epsilon)\hat{A}_t^{\text{GAE}}\right) \right]
\end{equation}
where $r_t(\theta) = \frac{\pi_\theta(a_{i,t}|o_{i,t})}{\pi_{\theta_{\text{old}}}(a_{i,t}|o_{i,t})}$ is the probability ratio and $\epsilon = 0.2$ is the clipping parameter.

A 2-stage curriculum is employed during training, splitting the total episode budget 60\%/40\%:
\begin{itemize}
    \item \textbf{Stage 1 (first 60\% of episodes)}: Spatial COLREGs pre-training under clean communication ($\lambda = 1.0$).
    \item \textbf{Stage 2 (remaining 40\%)}: Communication loss resilience fine-tuning under stochastic degradation ($\lambda \in [0.0, 1.0]$).
\end{itemize}

Table \ref{tab:hyperparams} details the vessel parameters and MAPPO hyperparameters used across all experiments.

\begin{table}[t]
\caption{ASV Hydrodynamic Coefficients \& MAPPO Hyperparameters}
\label{tab:hyperparams}
\centering
\begin{tabular}{llc}
\toprule
\textbf{Parameter Category} & \textbf{Parameter Symbol / Name} & \textbf{Value} \\
\midrule
\multirow{5}{*}{3-DOF MMG Dynamics} & Vessel Length ($L$) & 150.0 m \\
 & Displacement Mass ($m$) & $1.5 \times 10^7$ kg \\
 & Added Mass ($m_x, m_y$) & $5.0 \times 10^4$ kg, $1.0 \times 10^5$ kg \\
 & Max Propeller RPM ($n_{\text{max}}$) & 150 RPM \\
 & Max Rudder Angle ($\delta_{\text{max}}$) & $\pm 30.0^\circ$ \\
\midrule
\multirow{12}{*}{MAPPO / GAT Setup} & Discount Factor ($\gamma$) & 0.99 \\
 & GAE Parameter ($\lambda_{\text{gae}}$) & 0.95 \\
 & PPO Clip Ratio ($\epsilon$) & 0.20 \\
 & PPO Update Epochs & 4 \\
 & Batch Size & 1 full episode (on-policy, no minibatching) \\
 & Optimizer & Adam \\
 & Learning Rate & $3 \times 10^{-4}$ \\
 & Value Loss Weight & 0.5 \\
 & Entropy Coefficient & 0.01 \\
 & Actor/Critic Hidden Width & 64 \\
 & GAT Attention Heads ($K$) & 4 \\
 & Binary Payload Size & 128 bits \\
\bottomrule
\end{tabular}
\end{table}

\section{Experimental Evaluation}
\label{sec:experimental-evaluation}
Unless noted otherwise, all results in this section use a 2-vessel head-on encounter scenario ($N=2$ in the POMDP formulation of Section~\ref{sec:mmg-dynamics}): both vessels start 150m from the encounter point on a reciprocal course and must execute a COLREGs-compliant avoidance maneuver. This spacing has been retuned twice, in opposite directions, as the vessel dynamics of Section~\ref{sec:mmg-dynamics} changed: an original 4000m separation was tightened to 1600m when a first dynamics correction made the yaw response physically realistic but still slow, then tightened further to the current 300m gap after the IMO turning-circle fix (Section~\ref{sec:limitations}) made that yaw response an order of magnitude faster, which let every method clear enough distance to saturate the safety score at the earlier 1600m gap; see Section~\ref{sec:limitations} for the full account.

Every comparison in this section is reported in terms of one safety score,
\begin{equation}
J(\lambda) = \text{clip}\!\left(\frac{1}{|S|}\sum_{s \in S} d^{\min}_s(\lambda) \Big/ 1000\text{m},\; 0,\; 1\right),
\label{eq:safety_score}
\end{equation}
where $\lambda \in [0,1]$ is the communication-quality level a given evaluation episode is run under (Section~\ref{sec:mmg-dynamics}'s $p_{\text{msg}}$/comms channel), $S$ is the set of evaluation episodes (seeds) run at that level, and $d^{\min}_s(\lambda)$ is episode $s$'s true minimum pairwise vessel separation over its full length -- the real Euclidean distance between vessels at its single closest point, not a per-step projected closest-point-of-approach. $J$ is normalized against a 1000m reference and clipped to $[0,1]$; every reported score below is at least one method's episodes coming no closer than 1000m together (a fully safe pass), the floor is two vessels that collide or nearly do. $J$ is computed post-hoc from a completed episode's real trajectory and never appears inside any policy's reward (Eq.~\ref{eq:reward}), so no method in this paper is directly optimized against the metric it is evaluated on.

\subsection{Training Performance \& Convergence}
\label{sec:training-budget}
An earlier version of this evaluation trained every checkpoint for 150 episodes across 4 seeds, a budget chosen for computational tractability rather than validated as sufficient, with a separate 200-episode/3-seed run purely to characterize convergence behavior (Fig.~\ref{fig:extended_training}). Investigating why Rule-Based COLREGs reproducibly outperformed every RL-trained method at that budget (Section~\ref{sec:limitations}) surfaced a curriculum-training bug: \texttt{TwoStageCurriculumTrainer}'s Stage 2 fine-tuning phase reset the environment with the identical seed on every one of its episodes, so 40\% of every RL policy's training silently replayed one episode regardless of the requested episode count. We fixed this (a \texttt{seed\_offset} threaded through both training stages) and re-ran the comparison at two budgets: the original 150 episodes, now with the fix applied (\emph{e150}), and a realistic $3.3\times$ increase to 500 episodes (\emph{e500}), both across 5 independent seeds (42/100/200/300/400). Table~\ref{tab:budget_table} reports the results: at e150, every method remains within a similar range to the original (uncorrected) comparison, including Rule-Based COLREGs; at e500, MARLIN-Twin and every other RL-trained method's safety score rises to $0.089$--$0.097$, decisively above Rule-Based COLREGs' unchanged $0.085$ (Rule-Based has no learned parameters and is not affected by training budget). Section~\ref{sec:resilience-benchmark} reports the statistical test confirming this separation is significant, not an artifact of the larger budget alone happening to favor one run.

\begin{table}[h]
\caption{Safety Score $J(1.0)$ at the Pre-Fix (e150) and Post-Fix Realistic (e500) Training Budgets (Mean $\pm$ Std over 5 Training Seeds)}
\label{tab:budget_table}
\centering
\begin{tabular}{lcc}
\toprule
\textbf{Variant Model} & \textbf{e150 (150 eps)} & \textbf{e500 (500 eps)} \\
\midrule
\textbf{MARLIN-Twin (Full)} & 0.0838 $\pm$ 0.0237 & \textbf{0.0969 $\pm$ 0.0018} \\
Ablation 1 (Mean-Pool) & 0.1049 $\pm$ 0.0113 & 0.0972 $\pm$ 0.0025 \\
Ablation 2 (Flat MLP) & 0.0490 $\pm$ 0.0380 & 0.1021 $\pm$ 0.0090 \\
Ablation 3 (No Twin) & 0.0837 $\pm$ 0.0236 & 0.0969 $\pm$ 0.0018 \\
Independent PPO & 0.0624 $\pm$ 0.0352 & 0.0956 $\pm$ 0.0023 \\
MADDPG & 0.0893 $\pm$ 0.0131 & 0.0893 $\pm$ 0.0131 \\
MASAC (new, Section~\ref{sec:sac}) & --- & 0.0990 $\pm$ 0.0531 \\
Rule-Based COLREGs & 0.0851 $\pm$ 0.0000 & 0.0851 $\pm$ 0.0000 \\
\bottomrule
\end{tabular}
\end{table}

The convergence run of Fig.~\ref{fig:extended_training} (3 seeds, 200 episodes, purely diagnostic and not one of the compared checkpoints above) shows measurable improvement during Stage 1 pre-training and continued adaptation through the Stage 2 curriculum transition at episode 120 (the corresponding transition for the 150-episode budget elsewhere falls at episode 90, both at the curriculum's fixed 60\% split point); this run predates the Stage-2 seed-collapse fix above; and its qualitative shape -- improvement that continues past the original 150-episode cutoff -- is consistent with what the e150-vs-e500 comparison in Table~\ref{tab:budget_table} subsequently confirmed quantitatively.

\begin{figure}[t]
\centering
\includegraphics[width=0.48\textwidth]{fig10_extended_training_5k_seeds_ieee.pdf}
\caption{3-Seed Extended Curriculum MAPPO Training Curves with $\pm 1$ Standard Deviation Error Bands.}
\label{fig:extended_training}
\end{figure}

\subsection{Baseline Algorithm Comparison \& Resilience Index}
\label{sec:resilience-benchmark}
We benchmark MARLIN-Twin against IPPO \cite{yu2023mappo}, MADDPG \cite{lowe2017multi}, multi-agent SAC \cite{haarnoja2018soft} (Section~\ref{sec:sac}), and Rule-Based COLREGs \cite{tang2023colregs}, all at the e500 budget of Section~\ref{sec:training-budget}. We define a Coordination Resilience Index $R_{\text{resilience}}$, the ratio of the safety score's mean over $\lambda \in [0,1]$ to its value at full communication:
\begin{equation}
R_{\text{resilience}} = \frac{\int_{0}^{1} J(\lambda) \, d\lambda}{J(1.0)}
\end{equation}

\begin{figure}[t]
\centering
\includegraphics[width=0.48\textwidth]{fig9_benchmark_resilience_ieee.pdf}
\caption{Baseline Algorithm Coordination Resilience Comparison across MARLIN-Twin, Independent PPO, MADDPG, MASAC, and Rule-Based COLREGs, at the e500 training budget.}
\label{fig:resilience_benchmark}
\end{figure}

As shown in Fig. \ref{fig:resilience_benchmark}, the safety score $J(\lambda)$ now separates every RL-trained method from Rule-Based COLREGs, but does not reliably separate the four RL-trained methods from each other. Each trained method is evaluated across 5 independently trained seeds (42/100/200/300/400) and 10 evaluation seeds at full communication ($\lambda=1.0$): MARLIN-Twin ($0.0969 \pm 0.0018$), Independent PPO ($0.0956 \pm 0.0023$), MADDPG ($0.0893 \pm 0.0131$), and MASAC ($0.0990 \pm 0.0531$) all sit above Rule-Based COLREGs ($0.0851$, zero variance since it has no learned parameters). A Welch's $t$-test on MARLIN-Twin's 5 per-seed values against Rule-Based COLREGs' is significant ($p=0.00018$, Cohen's $d=8.5$, a very large effect); the same test against MADDPG is not ($p=0.31$, $d=0.72$), and against SAC is not either ($p=0.63$, $d=0.33$, SAC's cross-seed variance is far larger than any other method's, discussed in Section~\ref{sec:sac}) -- so while MARLIN-Twin now reliably beats the deterministic baseline, it remains statistically indistinguishable from the other learned methods, exactly as at the earlier budget. This is the central finding of this evaluation's second pass, and it reverses the first pass's headline result. That first pass, after fixing a missing per-seed re-initialization bug, found Rule-Based COLREGs ($0.1497$) reproducibly above all three RL-trained methods then tested ($0.080$--$0.091$, at the 150-episode/4-seed budget) -- we read that result as the more trustworthy one specifically because it was less flattering to our own method, and reported it as such. Investigating why a rule-based controller so reliably outperformed reward-shaped RL surfaced the curriculum Stage-2 seed-collapse bug of Section~\ref{sec:training-budget}: 40\% of every RL policy's training had silently replayed one episode regardless of budget. After fixing that bug and retraining at a realistic 500-episode budget, RL's disadvantage does not merely shrink -- it inverts, with MARLIN-Twin now significantly ahead. We disclose both reversals, not only the second, because a result that changes twice under closer scrutiny is worth reporting as it happened, not smoothed into a single final narrative.

The Coordination Resilience Index $R_{\text{resilience}}$ remains $1.0000 \pm 0.0000$ for every trained (RL) method and $0.9701$ for Rule-Based COLREGs at this budget -- none of the four RL policies show any measurable change in safety score across the six communication-degradation levels tested, only Rule-Based COLREGs does (a small but genuine variation). This invariance persists across both training budgets and both this and the previous evaluation pass; Section~\ref{sec:limitations} discusses why it should be read as a property of this 2-vessel head-on scenario and these policies' learned behavior, not as evidence the underlying mechanism cannot show resilience differentiation -- Section~\ref{sec:harder-scenarios}'s harder-scenario battery was designed in part to test that directly.

\subsection{Multi-Agent SAC: Testing a Fifth, Genuinely Different Algorithm}
\label{sec:sac}
Sections~\ref{sec:training-budget}--\ref{sec:resilience-benchmark} left one question open: three algorithms (MAPPO/MARLIN-Twin, MADDPG, Independent PPO) landing in the same statistically-indistinguishable band, even after the training-budget fix, is at least as consistent with a confound shared by all three (which the budget fix targeted) as with ``RL genuinely cannot exceed MADDPG/IPPO here.'' To test whether a fourth, structurally different algorithm changes this picture, we add multi-agent Soft Actor-Critic (SAC) \cite{haarnoja2018soft}: a maximum-entropy, off-policy method with a stochastic squashed-Gaussian actor, twin centralized critics, and an automatically tuned entropy coefficient, sharing MADDPG's replay-buffer and centralized-critic infrastructure but requiring neither a target actor nor a target encoder (next actions are resampled from the current stochastic actor rather than a lagged target network).

At the e500 budget, MASAC scores $0.0990 \pm 0.0531$ (Table~\ref{tab:budget_table}) -- nominally the highest mean of any method, but with by far the largest cross-seed standard deviation, more than $4\times$ MARLIN-Twin's. This variance is large enough that MASAC is not statistically distinguishable from Rule-Based COLREGs either ($p=0.63$), despite its mean sitting well above it: a wide per-seed spread can leave a method simultaneously ``numerically highest'' and ``not reliably better than the deterministic baseline,'' and both descriptions are true here. We attribute the variance to SAC's automatically-tuned entropy coefficient, which in preliminary runs occasionally grew very large on individual seeds (observed values into the hundreds) rather than settling toward zero as training progressed -- a known failure mode of automatic entropy tuning when the action space has a partially inert region (here, the throttle channel's mapping saturates over roughly 20\% of its raw range; see \texttt{VesselAgentWrapper.build\_action} in the released code), which we did not correct in this study to keep the algorithm-vs-algorithm comparison on equal footing with the other methods' unmodified action interface. MASAC's addition confirms that a fourth, architecturally distinct algorithm does not change the qualitative picture: every RL method examined so far beats Rule-Based COLREGs at e500, and no pair of RL methods is reliably distinguishable from each other.

\subsection{Real AIS Trajectory Validation}
We evaluate the EKF/JPDA Digital Twin using a real NOAA MarineCadastre Automatic Identification System trajectory \cite{marinecadastre2017ais, huang2025ais}: 48 consecutive real position reports (2017-01-20, $\sim$61-71s apart) broadcast by a vessel underway at 8.8-9.8 knots, drawn unmodified from NOAA's public-domain \texttt{AIS\_2017\_01\_Zone01} bulk release, with a 300s communication blackout injected over the middle third of the run. Fig. \ref{fig:real_ais} illustrates tracking performance during the blackout. The EKF Digital Twin maintains an overall position RMSE of approximately 497m against the real recorded trajectory, rising to approximately 1539m during the blackout window itself, using the vessel's own most recent reported speed and heading to dead-reckon forward across the real (irregular) elapsed time between reports.

\begin{figure}[t]
\centering
\includegraphics[width=0.48\textwidth]{fig11_real_ais_validation_ieee.pdf}
\caption{Real NOAA MarineCadastre AIS Trajectory Validation during an Injected 300s Blackout Period ($\sim$497m overall EKF RMSE vs $\sim$1539m during the blackout window).}
\label{fig:real_ais}
\end{figure}

\begin{figure}[t]
\centering
\includegraphics[width=0.48\textwidth]{fig12_ablation_study_ieee.pdf}
\caption{4-Variant Model Ablation Study comparing MARLIN-Twin against Mean-Pooling GNN, Flat MLP, and No Digital Twin variants.}
\label{fig:ablation_study}
\end{figure}

\subsection{4-Variant Model Ablation Study}
To evaluate the contribution of each system component, we conduct a 4-variant ablation study (Fig. \ref{fig:ablation_study}):
\begin{enumerate}
    \item \textbf{MARLIN-Twin (Full Proposed)}: GAT + EKF Digital Twin.
    \item \textbf{Ablation 1 (Mean-Pooling GNN)}: Replaces GAT attention with uniform mean pooling.
    \item \textbf{Ablation 2 (Flat MLP Policy)}: Replaces GNN graph encoding with flat vector MLP.
    \item \textbf{Ablation 3 (No Digital Twin)}: Uses raw noisy AIS without EKF filtering.
\end{enumerate}

Table \ref{tab:ablation_table} summarizes the quantitative metrics across all variants at the e500 budget (Section~\ref{sec:training-budget}), each averaged over 5 independently trained seeds (42/100/200/300/400) rather than a single training run, with the across-seed standard deviation reported alongside each mean, at full communication ($\lambda=1.0$). One finding is robust; the others are not, and one reverses under the fixed pipeline relative to the earlier (e150) evaluation. First, Ablation 3 (No Digital Twin) remains not merely statistically similar to the full model but bit-for-bit identical at the weight level -- we re-verified this directly at e500, at every one of the 5 seeds, after the Section~\ref{sec:training-budget} curriculum fix. Disabling the Digital Twin EKF, both during training and evaluation, still produces zero measurable effect here, for the same reason as before: this scenario never exercises a communication-degraded episode during either training or a lambda sweep in which the twin's filtered estimate would change the policy's behavior. Second, Ablation 1 (Mean-Pooling, $0.0972 \pm 0.0025$) and the full model ($0.0969 \pm 0.0018$) are now statistically indistinguishable -- the clear separation the e150 evaluation reported (Mean-Pooling $0.1161 \pm 0.0082$ vs.\ full model $0.0796 \pm 0.0272$) does not survive the realistic training budget, so we no longer read GAT attention as costing safety margin relative to uniform pooling in this scenario. Third, Ablation 2 (Flat MLP, $0.1021 \pm 0.0090$) also converges to the same range at e500, reversing the e150 trend (where it scored below the full model). All four variants -- Full, Mean-Pool, Flat MLP, and No-Twin -- land within $0.097$--$0.102$ at e500, a narrower spread than any pair's combined standard deviation; this evaluation does not support ranking any of the four architectural choices against the others at the realistic training budget, in contrast to the e150 budget's (since-reversed) separations. $R_{\text{resilience}}$ is $1.0000 \pm 0.0000$ for Full, Mean-Pool, and No-Twin; Ablation 2 (Flat MLP) is $1.0003 \pm 0.0006$, close to invariant like the others.

\begin{table}[h]
\caption{Quantitative Performance Metrics of Ablation Variants at Full Communication, e500 Budget (Mean $\pm$ Std over 5 Training Seeds)}
\label{tab:ablation_table}
\centering
\begin{tabular}{lcc}
\toprule
\textbf{Variant Model} & \textbf{Safety $J(\lambda)$} & \textbf{$R_{\text{resilience}}$} \\
\midrule
\textbf{MARLIN-Twin (Full)} & \textbf{0.0969 $\pm$ 0.0018} & \textbf{1.0000 $\pm$ 0.0000} \\
Ablation 1 (Mean-Pool) & 0.0972 $\pm$ 0.0025 & 1.0000 $\pm$ 0.0000 \\
Ablation 2 (Flat MLP) & 0.1021 $\pm$ 0.0090 & 1.0003 $\pm$ 0.0006 \\
Ablation 3 (No Twin) & 0.0969 $\pm$ 0.0018 & 1.0000 $\pm$ 0.0000 \\
\bottomrule
\end{tabular}
\end{table}

\subsection{Multi-Vessel Scenario Generalization}
\label{sec:multivessel}
All results above use the 2-vessel head-on scenario of Section~\ref{sec:experimental-evaluation}. To test whether the e500-budget reversal of Section~\ref{sec:resilience-benchmark} generalizes beyond a single pairwise encounter, we retrain MARLIN-Twin (GAT), Independent PPO, and the flat-MLP ablation (Ablation 2) from scratch on a denser 4-vessel channel-convergence scenario -- two opposing traffic lanes with tighter longitudinal spacing, producing multiple simultaneous pairwise encounters rather than one isolated pair -- at the same e500 budget (500 episodes, 5 independent seeds: 42/100/200/300/400) and the same safety-score evaluation. An earlier version of this table, produced before the curriculum seed-collapse bug (Section~\ref{sec:training-budget}) was found, reported this check as inconclusive; we re-ran it here under the fixed pipeline.

Table~\ref{tab:multivessel_table} reports the results: MARLIN-Twin (GAT, $0.1497 \pm 0.0590$), Independent PPO ($0.1120 \pm 0.0569$), and the flat-MLP ablation ($0.1253 \pm 0.0682$) still have $\pm 1$ standard deviation ranges that overlap each other substantially, so this check still does not support ranking the three architectures against one another. It does, however, now show a real, non-overlapping separation from Rule-Based COLREGs in the opposite direction from Section~\ref{sec:resilience-benchmark}'s primary result: Rule-Based COLREGs ($0.2652$, zero variance) sits clearly above every RL method's full range here. Unlike the 2-vessel head-on scenario, this channel-convergence layout gives every vessel a give-way duty that points the same direction regardless of which other vessel triggers it -- Rule-Based COLREGs' single-nearest-neighbor decision rule (Section~\ref{sec:harder-scenarios}) is never asked to arbitrate between two \emph{different} simultaneous obligations here, so its structural limitation is not exercised. We read this as a genuine, scenario-specific finding, not a contradiction of Section~\ref{sec:resilience-benchmark}: Section~\ref{sec:harder-scenarios}'s harder-scenario battery was designed specifically to test whether a scenario that \emph{does} exercise that limitation reverses this outcome, and finds two that do. $R_{\text{resilience}}$ is $1.0000 \pm 0.0000$ for MARLIN-Twin and Independent PPO; the flat-MLP ablation is $1.0008 \pm 0.0011$, and Rule-Based COLREGs is $1.0002$ -- all close to invariant here as elsewhere.

\begin{table}[h]
\caption{Safety Score on the 4-Vessel Channel-Convergence Scenario, e500 Budget (Mean $\pm$ Std over 5 Training Seeds)}
\label{tab:multivessel_table}
\centering
\begin{tabular}{lcc}
\toprule
\textbf{Policy} & \textbf{Safety $J(\lambda)$} & \textbf{$R_{\text{resilience}}$} \\
\midrule
MARLIN-Twin (GAT) & 0.1497 $\pm$ 0.0590 & 1.0000 $\pm$ 0.0000 \\
Independent PPO & 0.1120 $\pm$ 0.0569 & 1.0000 $\pm$ 0.0000 \\
Ablation 2 (Flat MLP) & 0.1253 $\pm$ 0.0682 & 1.0008 $\pm$ 0.0011 \\
Rule-Based COLREGs & \textbf{0.2652 $\pm$ 0.0000} & 1.0002 $\pm$ 0.0000 \\
\bottomrule
\end{tabular}
\end{table}

\subsection{Reward Reweighting Study}
\label{sec:reward-reweighting}
The reward MARLIN-Twin is trained against, $r_i = w_{\text{safety}} r_i^{\text{safety}} + w_{\text{colregs}} r_i^{\text{colregs}} + w_{\text{eff}} r_i^{\text{eff}}$ (Eq.~\ref{eq:reward}), does not directly optimize $J(\lambda)$; $r_i^{\text{safety}} = -\exp(-d_{\text{cpa}}/200\text{m})$ is within $e^{-5}$ of zero for any projected closest-approach beyond $\sim$1000m and only meaningfully negative inside $\sim$500m, while $r_i^{\text{eff}} = -d_{\text{goal}}/5000\text{m}$ pays a flat $+1.0$ every step once the route's single waypoint is reached (at $\sim$300m and 8~m/s, around step 37 of a 500-step episode) -- so for most of an episode, efficiency dominates safety's gradient by construction. To test whether this asymmetry, rather than the algorithm, was suppressing safety performance, we sweep the three weights across 8 configurations, MARLIN-Twin only, 3 seeds, e500 budget (Table~\ref{tab:reward_table}): A0 is the published weighting $(2,1,1)$; A1--A5 vary the weights directly; A6--A7 instead redefine $r_i^{\text{safety}}$ to use the true per-vessel minimum pairwise separation rather than the projected closest-point-of-approach -- the same quantity $J(\lambda)$ measures, a change to the reward's \emph{definition}, kept in its own rows rather than folded into the pure weight sweep.

\begin{table}[h]
\caption{Reward-Reweighting Sweep, MARLIN-Twin, e500 Budget (Mean $\pm$ Std over 3 Seeds)}
\label{tab:reward_table}
\centering
\begin{tabular}{lccc}
\toprule
\textbf{Arm} & \textbf{Weights} & \textbf{$J(1.0)$} & \textbf{Completion} \\
\midrule
A0 (published) & $(2, 1, 1)$ & 0.0977 $\pm$ 0.0018 & 0\% \\
A1 & $(2, 1, 0)$ & 0.0984 $\pm$ 0.0015 & 0\% \\
A2 & $(2, 1, 0.1)$ & 0.0979 $\pm$ 0.0019 & 0\% \\
A3 & $(8, 1, 1)$ & 0.0762 $\pm$ 0.0143 & 0\% \\
A4 & $(2, 4, 1)$ & 0.0993 $\pm$ 0.0000 & 0\% \\
A5 & $(8, 1, 0.1)$ & 0.0512 $\pm$ 0.0349 & 33\% \\
A6 (true-sep.) & $(2, 1, 0.1)$ & 0.0983 $\pm$ 0.0010 & 0\% \\
A7 (true-sep.) & $(2, 1, 1)$ & 0.0986 $\pm$ 0.0009 & 0\% \\
\midrule
Rule-Based COLREGs & --- & 0.0851 $\pm$ 0.0000 & 0\% \\
\bottomrule
\end{tabular}
\end{table}

Every arm's safety score lands within $0.076$--$0.099$ except A3 and A5, both of which upweight the safety term to $8\times$ without a compensating fix elsewhere; A5, which additionally cuts the efficiency weight to $0.1$, is clearly unstable (per-seed $J(1.0)$ of $0.016$, $0.099$, $0.039$) rather than merely lower on average. Naively upweighting safety alone destabilizes training rather than reliably improving it. Every arm's route-completion rate is $0\%$ except A5's ($33\%$), but every arm -- including Rule-Based COLREGs -- shows the same $0\%$ rate on this single deterministic evaluation episode, and mean vessel speed stays in the $6.5$--$7.4$~m/s range for every RL arm (vs.\ Rule-Based COLREGs' $8.9$~m/s) with no arm showing the near-zero speed a degenerate ``loiter to inflate $J$'' solution would require; we read the uniform $0\%$ completion rate as an artifact of this strict single-episode 50m-radius metric (which Rule-Based COLREGs also fails) rather than evidence of gaming. A6/A7's true-separation reward performs statistically indistinguishably from A0/A2, neither confirming nor ruling out that the reward/metric mismatch identified above materially affects the result at this budget. Overall, this sweep finds the published weighting already performs near the best of the range tested, and that the safety-vs-rule-based gap closed in Section~\ref{sec:resilience-benchmark} is not attributable to an easy reward-weighting fix that the published configuration was simply missing.

\subsection{Harder Scenario Battery}
\label{sec:harder-scenarios}
Rule-Based COLREGs (Section~\ref{sec:framework}) reduces every neighbor to the single nearest vessel by Euclidean distance and reacts only to that one; it cannot represent, let alone resolve, two simultaneous conflicting COLREGs obligations (give-way to one vessel while stand-on to another). The 2-vessel head-on scenario and the 4-vessel channel-convergence scenario (Section~\ref{sec:multivessel}) both structurally avoid ever presenting this case -- channel convergence gives every vessel a duty pointing the same direction regardless of which neighbor triggers it. To test scenarios that do exercise this limitation, we screen Rule-Based COLREGs (which needs no training) plus an untrained-random-policy floor across five candidates -- the existing 4-vessel channel-convergence scenario, a 6-vessel port-approach scenario (three $120^\circ$ sectors converging on one point, so a vessel simultaneously has one give-way and one stand-on neighbor), a 6-vessel head-on scenario, and a restricted-visibility crossing scenario at $N=4$ and $N=6$ -- before committing any training compute (Table~\ref{tab:screen_table}).

\begin{table}[h]
\caption{Scenario Screen: Rule-Based COLREGs vs. an Untrained Random Policy}
\label{tab:screen_table}
\centering
\begin{tabular}{lccc}
\toprule
\textbf{Candidate} & \textbf{$N$} & \textbf{Rule-Based $J(1.0)$} & \textbf{Random $J(1.0)$} \\
\midrule
Channel convergence & 4 & 0.2583 & 0.0760 \\
Port approach & 6 & \textbf{0.1111} & \textbf{0.1867} \\
Head-on & 6 & 0.0091 & 0.0048 \\
Restricted vis.\ crossing & 4 & 0.2027 & 0.1926 \\
Restricted vis.\ crossing & 6 & 0.1887 & 0.1366 \\
\bottomrule
\end{tabular}
\end{table}

Two candidates stand out. The 6-vessel port-approach scenario is the most striking: Rule-Based COLREGs scores \emph{below} its own random-policy floor ($0.1111$ vs.\ $0.1867$) -- doing measurably worse than an untrained policy, the strongest possible evidence that its single-nearest-neighbor rule fails outright here. The 6-vessel restricted-visibility crossing scenario also shows Rule-Based COLREGs meaningfully above the random floor but well below the other candidates' absolute scores, indicating real but bounded difficulty. We escalate both to full training (5 variants -- MARLIN-Twin, Independent PPO, the flat-MLP ablation, MADDPG, and MASAC -- $\times$ 5 seeds $\times$ 500 episodes), rather than only the single scenario a stricter selection rule would have picked, since both are genuinely informative and the training compute to do so was affordable (Table~\ref{tab:battery_table}).

\begin{table}[h]
\caption{Harder-Scenario Escalation: Safety Score $J(1.0)$ (Mean $\pm$ Std over 5 Seeds)}
\label{tab:battery_table}
\centering
\begin{tabular}{lcc}
\toprule
\textbf{Variant} & \textbf{Port Approach ($N{=}6$)} & \textbf{Restr.\ Vis.\ Crossing ($N{=}6$)} \\
\midrule
MARLIN-Twin & 0.2288 $\pm$ 0.0143 & 0.2878 $\pm$ 0.1099 \\
Independent PPO & \textbf{0.3258 $\pm$ 0.0152} & \textbf{0.3800 $\pm$ 0.0191} \\
Flat MLP & 0.2515 $\pm$ 0.0188 & 0.3641 $\pm$ 0.0166 \\
MADDPG & 0.2244 $\pm$ 0.0295 & 0.3772 $\pm$ 0.0223 \\
MASAC & 0.1762 $\pm$ 0.0410 & 0.3014 $\pm$ 0.1478 \\
\midrule
Rule-Based COLREGs & 0.1316 $\pm$ 0.0000 & 0.1906 $\pm$ 0.0000 \\
\bottomrule
\end{tabular}
\end{table}

On both scenarios, \emph{every} trained RL variant clearly exceeds Rule-Based COLREGs, confirming the screen's prediction with genuinely trained policies rather than the untrained comparison alone. Independent PPO scores highest on both, not MARLIN-Twin -- we report this as found rather than only reporting the method this paper proposes; it is consistent with Section~\ref{sec:resilience-benchmark} and Section~\ref{sec:sac}'s finding that no RL method examined in this paper is reliably distinguishable from any other. Combined with Section~\ref{sec:multivessel}'s channel-convergence result, the pattern across all three multi-vessel scenarios tested is consistent: Rule-Based COLREGs wins on the one scenario (channel convergence) that never presents it with conflicting simultaneous obligations, and loses on both scenarios (port approach, restricted-visibility crossing) that do -- exactly the structural weakness this battery was designed to exercise, not an outcome engineered by picking scenarios until RL wins. Had every candidate instead shown Rule-Based COLREGs winning, that would have been an equally valid, reportable finding; it is not what we found.

\subsection{Limitations}
\label{sec:limitations}
Several results in this paper reflect current implementation limitations rather than validated claims, and should be read accordingly:

\begin{itemize}
    \item \textbf{Hydrodynamic coefficient realism.} The MMG added-mass/damping coefficients used throughout (Section~\ref{sec:mmg-dynamics}) remain hand-tuned engineering estimates for CARGO and USV -- the only two vessel types any experiment in this paper instantiates -- not derived from captive-model-test data or a documented regression method such as Clarke, Gedling, and Hine's (1983) hull-form-based estimation approach \cite{clarke1983manoeuvring}, Yoshimura et al.'s empirical MMG-parameter formulas \cite{wang2024mmg}, or Chame et al.'s (2025) more recent hull-form regression \cite{chame2025linearderivative}; this caveat is unchanged from our first evaluation attempt. What changed is IMO Res. MSC.137(76) compliance itself. We had previously traced a $\sim$7.6~km turning-circle tactical diameter to a yaw spin-up time constant an order of magnitude longer than the loop-closing time, and an exhaustive search across the model's existing damping/rudder-gain parameters and several added cross-coupling terms appeared to show every combination trading turning-circle compliance for zig-zag stability, with none satisfying both. That search, however, rested on two flawed measurement tools we subsequently found and fixed. First, \texttt{run\_zigzag\_test}'s phase 2/3 rudder-reversal detection compared a raw, non-modulo-wrapped heading difference against its $\pm 10^\circ$ threshold, which breaks discontinuously once a fast-turning candidate's true relative heading passes $\pm 180^\circ$ -- exactly the regime a vessel tight enough to meet the turning-circle criterion operates in. Second, our own full-trajectory stability check -- added specifically to catch coefficient sets that pass the two officially-checked overshoots while diverging afterward -- mistook the vessel's own held-rudder turning circle for divergence: the zig-zag test executes only the two IMO-specified reversals and then holds the rudder fixed for the remainder of the test window, so any vessel maneuverable enough to complete its turning circle within that window sweeps its heading through the full $\pm 180^\circ$ range every lap, which the check flagged as instability regardless of whether the response was actually diverging. After fixing both -- reusing one wrap-corrected relative heading throughout the reversal logic, and re-deriving stability from overshoot bounded across repeated reversals rather than raw heading extremes -- the coefficient candidates the original search had rejected as unstable were, in fact, stable: bounded overshoot (peak 1.84$^\circ$) across 10 consecutive zig-zag reversal cycles, with both official IMO numbers holding simultaneously (466~m tactical diameter, 10.0$^\circ$/19.6$^\circ$ overshoots; Section~\ref{sec:mmg-dynamics}). We adopted the corrected values for both vessel types used in every experiment in this paper (CARGO: yaw damping $N_r=-1\times10^9$ and rudder-moment gain $8000$, from $-2\times10^4$/$24.87$; USV: $N_r=-4\times10^7$ and gain $24000$, from $-4\times10^3$/$49.74$) and retrained every checkpoint under them, since the change materially affects how sharply a vessel can turn in response to a policy's rudder command. The negative evidence we previously reported for added $N_v\cdot v$/$Y_r\cdot r$ cross-coupling terms was not re-verified against the corrected test harness and should not be read as settled either way; reaching IMO compliance did not require them, so we do not revisit them here. The coefficients themselves remain a direct-simulation-search result, not a captive-model-test- or regression-derived one, so the hand-tuned-estimate caveat above still applies -- what this closes is IMO sea-trial compliance specifically, not hydrodynamic realism generally.
    \item \textbf{Resilience index interpretation.} $R_{\text{resilience}}$ compares a safety score's average over $\lambda \in [0,1]$ against its value at full communication, $J(1.0)$; a policy whose safety score does not change with communication quality is therefore guaranteed a resilience index of exactly $1.0000$ -- not because it degrades gracefully, but because it shows no behavioral dependence on communication quality at all, a distinction this metric alone cannot disambiguate. At the e500 budget (Section~\ref{sec:training-budget}), $R_{\text{resilience}} = 1.0000$ still holds for every one of the four RL-trained methods on the primary 2-vessel scenario (Section~\ref{sec:resilience-benchmark}), on the 4-vessel channel-convergence scenario (Section~\ref{sec:multivessel}), and on both harder scenarios of Section~\ref{sec:harder-scenarios} -- port approach and restricted-visibility crossing -- where only Rule-Based COLREGs shows any measurable variation (and even there, only a few points in the third decimal place). This is a stronger version of the same finding an earlier version of this evaluation reported at 150 episodes and attributed, tentatively, to insufficient training: the invariance now persists across two training budgets, five algorithms (including two, MADDPG and SAC, that were not part of the original finding), and four distinct scenarios, including two specifically designed to be harder and more ambiguous than the original. We no longer read "needs more training" or "needs a harder scenario" as plausible explanations; the pattern is consistent instead with every policy examined in this paper acting close to deterministically at evaluation time (all methods here select actions greedily, without exploration, when evaluated) in a way that does not condition materially on the noisier, perceived-neighbor-state input degraded communication produces, at least not enough to move the true minimum pairwise separation the safety score measures. Determining whether genuine resilience differentiation is achievable at all under this metric would likely require either evaluating with stochastic (exploratory) action selection retained, or a metric sensitive to smaller behavioral changes than $J(\lambda)$'s coarse separation-based reduction can detect. Fig.~\ref{fig:bandwidth_heatmap}'s flat team-reward curve is discussed where it appears (Section~\ref{sec:framework}) and is a related but distinct finding from this one, for the same underlying reason: its reward's own safety term was deliberately left uncorrected in this study, though Section~\ref{sec:reward-reweighting}'s A6/A7 arms test a corrected version without changing this section's conclusion.
    \item \textbf{A fourth round: a curriculum seed-collapse bug reversed the headline finding of the third round.} Investigating why Rule-Based COLREGs so reliably outperformed reward-shaped RL at the 150-episode budget (the one finding that survived the third round below) led us to \texttt{TwoStageCurriculumTrainer.train\_curriculum}, which runs Stage~2 (bandwidth-degradation fine-tuning) as a loop of single-episode calls to \texttt{MAPPOTrainer.train}. That method resets the environment with \texttt{seed=ep}, where \texttt{ep} is the loop index \emph{within} the current call -- always $0$ for a call given \texttt{n\_episodes=1} -- so every Stage~2 episode, 40\% of every RL policy's total training, silently reset to the identical seed regardless of how many episodes were requested; Stage~1 (60\% of training, run as one multi-episode call) was unaffected. We fixed this with a \texttt{seed\_offset} parameter threaded through both stages so Stage~1 and Stage~2 draw disjoint seed ranges, and reverified with a unit test that a 10-episode curriculum run now visits 10 distinct seeds (it previously visited 1). This bug is independent of the seeding bug the next paragraph describes -- that one affected \emph{initialization} across seeds; this one affected \emph{episode diversity within} a single seed's own Stage-2 training -- and affected every curriculum-trained variant (MARLIN-Twin and its three ablations, Independent PPO) but not MADDPG or MASAC, whose off-policy trainers do not use the curriculum wrapper. Retraining at the original 150-episode budget with this fix applied changes the earlier numbers only modestly (Table~\ref{tab:budget_table}'s e150 column); retraining at a realistic 500-episode budget changes them substantially, reversing Rule-Based COLREGs' advantage into a significant MARLIN-Twin advantage (Section~\ref{sec:resilience-benchmark}). Alongside this fix, we also corrected three smaller issues surfaced during the same investigation: Rule-Based COLREGs' action output was being passed through the same tanh-squashed-to-physical-units remapping every learned policy's raw network output requires, double-transforming its already-physical intended commands (this happened to not affect this paper's own published figures, which each separately bypassed the shared remapping for Rule-Based COLREGs, but would have silently affected any new evaluation script reusing the standard path, so we unified the convention rather than leave the inconsistency for a future script to inherit); two internal evaluators (used for training-time logging, not any figure in this paper) fed the same projected closest-point-of-approach the reward uses, rather than the true minimum pairwise separation Eq.~\ref{eq:safety_score} defines, into the safety-score formula, which we reconciled onto the same true-separation convention throughout; and the \texttt{overtaking} scenario generator assigned both vessels identical speed, which \texttt{classify\_encounter}'s speed-ordering test can never classify as a genuine overtaking encounter, so we gave the trailing vessel a materially higher speed and added a minimal Rule 13 (overtaking) response to Rule-Based COLREGs, which previously had no branch for that encounter type at all.
    \item \textbf{Retuning the encounter geometry for the turning-circle fix, and an earlier seeding bug, together drove three rounds of revised comparative rankings before the fourth round above.} The IMO turning-circle fix (previous bullet) raised CARGO/USV's yaw damping and rudder-moment gain by several orders of magnitude, so both vessel types can begin an avoidance turn almost immediately. At the head-on encounter geometry used throughout an earlier version of this evaluation (800m/side, 1600m gap), every policy -- including Rule-Based COLREGs, which has no learned parameters -- could clear more than 1000m of separation, saturating the safety score's 1000m-normalized ceiling for every method simultaneously (confirmed directly: 1285m closest approach for Rule-Based COLREGs, 1332m for MARLIN-Twin at that gap under the corrected dynamics) and erasing its ability to differentiate them. We retuned the gap to 150m/side (300m total) by direct simulation search, choosing the tightest value still consistent with a $\sim$100-150m closest approach for a reasonable avoidance policy, and retrained every affected checkpoint under it (the 4-vessel channel-convergence scenario's own geometry was separately verified not to need retuning, a $\sim$283m Rule-Based closest approach there already, under the same corrected dynamics). Retraining after this retightening produced, in an intermediate version of this evaluation, three seemingly clean-cut rankings: Rule-Based COLREGs scoring highest of all four baseline-comparison methods, the No-Digital-Twin ablation scoring higher than the full model, and the flat-MLP ablation scoring roughly $6.4\times$ higher than MARLIN-Twin's graph-attention policy at $N=4$ -- each with a reported cross-seed standard deviation of exactly (or very nearly) zero. That near-zero variance was the tell: investigating why 4 independently trained seeds should agree that closely surfaced a real bug -- \texttt{run\_retrain\_all\_baselines.py} and \texttt{run\_multivessel\_generalization.py} never called \texttt{seed\_everything(seed)} before constructing a policy, so the \texttt{seed} argument only ever named the checkpoint file; for 5 of the 6 baseline/ablation policy types, all 4 "independently trained" checkpoints were, in fact, bit-for-bit identical weights (confirmed by direct tensor comparison), and the sixth (MARLIN-Twin's GATPolicy) only showed real variance because its larger initialization happened to consume enough incidental global-RNG draws under our multiprocessing training pool to diverge across workers -- not because it, alone, was seeded correctly. We fixed this by adding the missing \texttt{seed\_everything(seed)} calls, verified directly that checkpoints now differ across seeds for every variant, and retrained and re-evaluated everything a third time. Two of the three intermediate rankings did not survive genuine seed variance and are now statistically inconclusive (Section~\ref{sec:experimental-evaluation}, Table~\ref{tab:multivessel_table}); the No-Digital-Twin result changed in an interesting way rather than simply vanishing -- it is not just statistically similar to the full model, but produces bit-for-bit identical trained weights, because this scenario never presents training or evaluation conditions under which the twin's filtering would change the policy's behavior (Table~\ref{tab:ablation_table}). Rule-Based COLREGs' advantage over all three RL methods is the one finding from the intermediate version that held up under correctly randomized seeds. We report all of this, including our own error and how we found it, because a paper's statistical claims are only as trustworthy as the seeding infrastructure producing them, and we would rather disclose that infrastructure failed once than have a reader discover it independently.
    \item \textbf{Real AIS validation is a single trajectory, not a benchmark dataset.} Fig.~\ref{fig:real_ais} now uses genuine NOAA MarineCadastre telemetry rather than a simulated trajectory, but it is one vessel's 51-minute excerpt with one injected blackout, not a multi-vessel or multi-condition real-world benchmark; the 497m/1539m RMSE figures characterize the EKF/JPDA estimator on this one real trajectory, not a general real-world accuracy claim. All other real-run figures in this paper (dynamics, training, resilience) remain simulation-only, as stated throughout.
    \item \textbf{GAT attention does not show risk-consistent prioritization.} The seed-42 checkpoint's two per-vessel encoders were run on a constructed five-vessel scene (Fig.~\ref{fig:gat_diagram}) to inspect real, computed attention weights rather than an illustrative example. One encoder (vessel 0) assigns its highest weight to the head-on neighbor; the other (vessel 1) assigns its highest weight to the overtaking neighbor instead. Since this checkpoint was trained for only 150 episodes on a 2-vessel scenario and is being evaluated here out-of-distribution on 5 vessels, this is consistent with the attention mechanism not yet having learned a \emph{stable, vessel-independent} risk-based prioritization, rather than evidence against the mechanism's capacity to do so with more training.
\end{itemize}

\section{Conclusion \& Future Work}
This paper presented MARLIN-Twin, a multi-agent reinforcement learning framework for autonomous vessel collision avoidance under communication outages, combining an EKF/JPDA Digital Twin with a Multi-Head Graph Attention Network and CTDE MAPPO optimization, and used it as a testbed to evaluate rather than assume how well this design delivers on its aims. We diagnosed and fixed a genuine turning-circle non-compliance with IMO Res. MSC.137(76), tracing it to a latent measurement bug in our own zig-zag test harness rather than a fundamental modeling limit, and a missing per-seed re-initialization bug that had left most of an earlier evaluation pass's "independently trained" checkpoints identical across seeds. That pass, corrected for both bugs, found Rule-Based COLREGs reproducibly outperforming every 150-episode-trained RL policy -- a result we reported as more trustworthy specifically because it was less flattering to our own method. Investigating \emph{that} result surfaced a fourth bug, independent of the first two: a curriculum-training seed-collapse fault that silently replayed one training episode for 40\% of every RL policy's budget regardless of how many episodes were requested (Section~\ref{sec:limitations}). Fixing it and retraining at a realistic 500-episode budget reverses the finding a second time: MARLIN-Twin now significantly outperforms Rule-Based COLREGs (safety score $0.097$ vs.\ $0.085$, $p=0.00018$, Cohen's $d=8.5$), and the same reversal holds for every RL method tested, including a fifth algorithm, multi-agent SAC, added specifically to test whether the reversal was architecture-specific. MARLIN-Twin remains statistically indistinguishable from MADDPG, Independent PPO, and SAC, and the three ablation variants (mean-pooling attention, a flat-MLP policy, and a disabled Digital Twin) that separated cleanly at 150 episodes converge to the same range at 500 -- so which architectural choice is best, among those tested, remains open even though whether RL beats this deterministic baseline is now resolved, on this scenario, in RL's favor.

We tested this reversal's robustness two further ways. First, sweeping the reward's safety/COLREGs/efficiency weights across eight configurations (Section~\ref{sec:reward-reweighting}) found the published weighting already near the best of the range tested, and that naively upweighting the safety term alone destabilizes training rather than improving it -- the reversal is not attributable to an easy reward-tuning fix the published configuration happened to miss. Second, we screened five harder, more ambiguous multi-vessel scenarios against a structural weakness we confirmed directly in Rule-Based COLREGs' own decision logic -- it collapses every neighbor to the single nearest vessel and cannot represent simultaneous conflicting COLREGs obligations -- and escalated the two most discriminative candidates to full training (Section~\ref{sec:harder-scenarios}). On both a converging multi-sector port-approach scenario and a restricted-visibility crossing scenario, every trained RL variant, not MARLIN-Twin alone, clearly outperforms Rule-Based COLREGs; a fourth multi-vessel scenario (channel convergence, Section~\ref{sec:multivessel}), which never presents conflicting simultaneous obligations, continues to favor Rule-Based COLREGs. We read this pattern -- Rule-Based COLREGs winning only on the one scenario structurally suited to its single-neighbor decision rule, and losing everywhere that rule is genuinely challenged -- as a coherent, if scenario-dependent, positive result for learned coordination generally, not for any one architecture specifically. The Coordination Resilience Index remains exactly invariant ($R_{\text{resilience}}=1.0000$) for every RL method across every scenario and both training budgets tested, including the two new harder scenarios explicitly designed to make it otherwise; we no longer read this as an artifact of insufficient training or an insufficiently hard scenario (Section~\ref{sec:limitations} discusses what would be needed to actually distinguish it).

Validation against a real NOAA MarineCadastre AIS trajectory demonstrates the EKF/JPDA Digital Twin's dead-reckoning fallback functions correctly during a blackout, with position error growing to roughly 1539m during the outage window. We consider disclosing every seeding, measurement-harness, and metric-consistency bug found across both evaluation passes -- and the fact that this paper's headline finding reversed twice as a direct result -- as much a contribution of this work as any single result: an evaluation's comparative claims are only as trustworthy as the randomization and measurement infrastructure producing them, and we would rather show that infrastructure failing and being caught than report only the version of events that looks cleanest in hindsight. Future work will extend the harder-scenario battery to a wider range of multi-vessel COLREGs geometries built from scratch (mixed crossing/overtaking/head-on encounters within one scene, rather than the two existing-scenario candidates escalated here), run the reward-reweighting sweep across the other algorithms tested rather than MARLIN-Twin alone, investigate why the Coordination Resilience Index remains exactly invariant under stochastic (exploratory) action selection rather than only the deterministic evaluation used throughout this paper, restructure the MMG force law to the standard linear-plus-cross-flow form so that published hull-form regressions (e.g. Clarke et al. \cite{clarke1983manoeuvring}; Yoshimura et al. \cite{wang2024mmg}; Chame et al. \cite{chame2025linearderivative}) can be substituted for the current hand-tuned damping coefficients, integrate 3D wave disturbances, and conduct physical ASV field trials.

\bibliographystyle{IEEEtran}
\bibliography{references}

\end{document}
